\Needspace{8\baselineskip}\originalchapter{分布函数、矩方程与闭合}\label{ch:moments}
\originalsection{相空间分布}
$f_s(\mathbf r,\mathbf v,t)\dd^3r\dd^3v$ 是相空间小体积的粒子数, $n_s=\int f_s\dd^3v$, 流速 $\mathbf u_s=n_s^{-1}\int\mathbf vf_s\dd^3v$. 以热速度 $\mathbf w=\mathbf v-\mathbf u_s$ 定义压力张量 $P_{ij}=m_s\int w_iw_jf_s\dd^3v$, 热流 $\mathbf q_s=(m_s/2)\int w^2\mathbf wf_s\dd^3v$. 分布不一定 Maxwell, 压力不一定标量.

粒子守恒的相空间动力学写为
\begin{equation}
 \partial_tf_s+\mathbf v\cdot\nabla f_s+\frac{q_s}{m_s}(\mathbf E+\mathbf v\times\mathbf B)\cdot\nabla_vf_s=C_s[f].
\end{equation}
无碰撞时 $C_s=0$ 为 Vlasov 方程, 沿特征线 f保持. ``无碰撞''不等于``没有相互作用'', 场依各分布产生并耦合回粒子.

\originalsection{从动力学到流体}
对上式积速度, 分布在速度无穷远足够衰减使分部积分边界零. 得连续性 $\partial_tn+\nabla\cdot(n\mathbf u)=0$. 乘 $m\mathbf v$ 再积, 把二阶矩拆成 $mn\mathbf u\mathbf u+\mathbf P$, 得
\begin{equation}
 mn(\partial_t+\mathbf u\cdot\nabla)\mathbf u
 =qn(\mathbf E+\mathbf u\times\mathbf B)-\nabla\cdot\mathbf P+\mathbf R.
\end{equation}
$\mathbf R=m\int\mathbf vC\dd^3v$ 是碰撞动量交换. 对整个封闭多物种体系 $\sum_s\mathbf R_s=0$, 但各物种可不同.

二阶矩演化含三阶热流, 三阶又含四阶, 因而矩阶层不自行闭合. 强局域碰撞和慢变化可支持各向同性 Maxwell 与标量 p. 绝热无热传导闭合 $D(pn^{-\gamma})/Dt=0$, 三维单原子 $\gamma=5/3$; 一维受约束热运动可取其他有效值. 电子等离子体波高相速闭合常给 $\gamma_e=3$, 不应混同静态等温 Boltzmann响应的 $\gamma_e=1$.

\originalsection{各向异性与双绝热}
强磁化、碰撞慢时, $\mathbf P=p_\perp(I-\mathbf b\mathbf b)+p_\parallel\mathbf b\mathbf b$. 单粒子磁矩及沿磁通管的体积变化, 在忽略热流的 CGL 闭合中给
\begin{equation}
 \frac{D}{Dt}\frac{p_\perp}{nB}=0,\qquad
 \frac{D}{Dt}\frac{p_\parallel B^2}{n^3}=0.
\end{equation}
第一式是横动能随 B 改变, 第二式含平行压缩. 它们有适用假设, 非任意碰撞less体系必然满足. 大的压力各向异性可能驱动 mirror 或 firehose, 不能用各向同性 MHD 的稳定结论替代.

\originalsection{两流体与广义 Ohm 定律}
电子和离子各写动量方程, 用准中性 $n_e=Zn_i$、电流 $\mathbf J=en_iZ(\mathbf u_i-\mathbf u_e)$, 在 $m_e\ll m_i$ 近似取质心速度 $\mathbf V\simeq\mathbf u_i$. 电子方程改写为
\begin{equation}\label{eq:ohm-general}
 \mathbf E+\mathbf V\times\mathbf B=\eta\mathbf J+\frac{\mathbf J\times\mathbf B}{en_e}
 -\frac{\nabla p_e}{en_e}+\frac{m_e}{n_ee^2}\frac{D\mathbf J}{Dt}+\cdots.
\end{equation}
各项分别为电阻、Hall、电子压强与电子惯性, 省略号含密度与对流修正. 哪些项可忽略须比较具体尺度, 不能仅因等离子体很热就宣布理想 Ohm. Hall 在离子皮肤深度尺度可能显著, 电子惯性在更小尺度显著; 两者是连接 MHD 与波、重联动力学的桥梁.
